% Detailed implementation note for the AmpliCol gen23 phase-space map.
% Compile with: pdflatex gen23_phase_space.tex
\documentclass[11pt,a4paper]{article}

\usepackage[T1]{fontenc}
\usepackage[utf8]{inputenc}
\usepackage{amsmath,amssymb,mathtools}
\usepackage{booktabs,array,geometry,hyperref}
\usepackage{xcolor}
\geometry{margin=27mm}
\hypersetup{colorlinks=true,linkcolor=blue!55!black,citecolor=blue!55!black,urlcolor=blue!55!black}
\setlength{\parskip}{0.55em}
\setlength{\parindent}{0pt}
\allowdisplaybreaks

\newcommand{\shat}{\hat{s}}
\newcommand{\Ph}{\mathrm{d}\Phi}
\newcommand{\dd}{\mathrm{d}}
\newcommand{\cM}{\mathcal{M}}
\newcommand{\code}[1]{\texttt{#1}}
\newcommand{\lam}{\lambda}
\newcommand{\ET}{E_{\mathrm{T}}}
\newcommand{\pt}{p_{\mathrm{T}}}
\newcommand{\mc}[1]{\mathcal{#1}}

\title{The \code{gen23} Phase--Space Generator\\
       Analytic derivation, implementation map, bounds, and audit}
\author{AmpliCol implementation note}
\date{\today}

\begin{document}
\maketitle

\begin{abstract}
This note documents the phase--space parametrisation implemented in the
\code{gen23} source.  The generator is a recursive,
colour-order-aware construction for massless incoming partons and an arbitrary
number of massive or massless final-state particles.  Its distinctive building
block is the invariant-variable $2\to3$ construction of Byckling and Kajantie;
ordinary $s$- and $t$-channel two-body blocks and a special double-$t$ block
are used where they are more natural.  All coordinate maps, Jacobian factors,
physical boundaries, cut-derived pre-bounds, and the inverse map are described.

The final sections deliberately separate exact kinematics from bounds that are
only conservative implementation aids.  They also record several concrete
discrepancies between the analytic acceptance region and the present code,
including one all-final-state lower bound that can remove accepted partonic
energies when pair-mass cuts are enabled.
\end{abstract}

\tableofcontents

\section{Scope and source of truth}

This document describes the current implementation in
\code{PhaseSpace/phase\_space\_gen23.f03}, rather than an idealised generic
``gen23'' algorithm.  In particular, a statement such as ``the generator uses
the Byckling--Kajantie variables'' is not enough to reproduce an AmpliCol
weight: the code also selects a topology from a colour ordering, applies
importance maps, inserts cut-based restrictions, chooses between a two-fold
azimuthal ambiguity, divides out powers of $2\pi$, and includes the partonic
flux factor.

The analytic invariant construction is based on
Byckling and Kajantie~\cite{BK}.  The code comments refer especially to their
Eqs.~(8), (10), (11), and (13).  Equation numbering in this note is local;
the displayed expressions are written in the variable convention actually
passed to the Fortran routines.

The map has the following intended domain.

\begin{itemize}
\item There are two \emph{massless} incoming particles, labelled $a$ and $b$
  (external labels 1 and 2 in the phase-space ordering), and $N$ final
  particles.  The code aborts for non-zero incoming masses.
\item $N\geq2$.  At fixed partonic energy the dimension is $3N-4$; a $2\to1$
  reaction needs a separate treatment and is not represented by this map.
\item A cyclic colour ordering contains each external label exactly once and
  contains labels 1 and 2.  It determines the recursive topology; it is not a
  mere presentation convention.
\item Final-state masses are supported.  Some cut pre-bounds are exact only
  for massless final legs; those qualifications are made explicit below.
\end{itemize}

\section{Lorentz-invariant phase space and normalisation}

We use metric $g=(+,-,-,-)$ and write
\begin{equation}
 p\mathbin{\cdot}q=p^0q^0-\boldsymbol p\mathbin{\cdot}\boldsymbol q,
 \qquad
 \lam(x,y,z)=x^2+y^2+z^2-2xy-2xz-2yz .
 \label{eq:lambda}
\end{equation}
The function in Eq.~\eqref{eq:lambda} is implemented by \code{lambda()}.
For a partonic process $a+b\to1+\cdots+N$, the standard Lorentz-invariant
measure is
\begin{equation}
 \dd\Phi_N(P;p_1,\ldots,p_N)=
 (2\pi)^4\delta^{(4)}\!\left(P-\sum_{r=1}^{N}p_r\right)
 \prod_{r=1}^{N}\frac{\dd^3\boldsymbol p_r}{(2\pi)^3 2E_r},
 \qquad P=p_a+p_b,
 \label{eq:lips}
\end{equation}
where $P^2=\shat$.  The partonic differential cross section used by the code is
\begin{equation}
 \dd\hat\sigma=\frac{1}{2\shat}\,
 \overline{|\cM|^2}\,\dd\Phi_N .
 \label{eq:flux}
\end{equation}
Thus the final division by \code{2d0*sqrtshat**2} in the source is the flux
factor, not part of the phase-space volume itself.

At fixed $\shat$, resolving four momentum-conservation delta functions leaves
\begin{equation}
 D_{\rm PS}=3N-4=3(n-2)-4,
 \label{eq:dimension}
\end{equation}
where \code{n} is the total number of external particles.  This is exactly
\code{this\%ndim} without PDFs.  With PDFs, the map additionally integrates
two Bjorken fractions and has $D_{\rm PS}+2$ adaptive coordinates.

It is useful to define a reduced measure in which the global $2\pi$ factors
are temporarily suppressed.  Every complete fixed-$\shat$ map constructed by
the source obeys
\begin{equation}
 \dd\Phi_N=\frac{\dd\overline\Phi_N}{(2\pi)^{3N-4}}.
 \label{eq:reducedmeasure}
\end{equation}
The source first accumulates $\dd\overline\Phi_N$ and, at the end of
\code{generate\_momenta}, divides the Jacobian by
\code{(2d0*pi)**(3*(next-2)-4)}.  Equations below quote reduced Jacobians when
they are meant to match an individual kernel directly.

\section{Masks, crossed sums, and colour-order topology}

\subsection{Bit-mask representation}

An integer mask $A$ labels a subset of external legs.  Bit $r-1$ represents
external label $r$.  Hence the singleton $r$ is \code{ibset(0,r-1)}, disjoint
masks are combined by integer addition, and the complement is
\begin{equation}
 \bar A=(2^n-1)-A .
\end{equation}
The arrays \code{pp(:,A)} and \code{invm(A)} respectively hold an intermediate
four-vector and its square.  For the invariant bookkeeping, the natural
crossed sum is
\begin{equation}
 P_A=\sum_{r\in A\cap\mc F}p_r-
     \sum_{c\in A\cap\{a,b\}}p_c,
 \qquad M_A^2=P_A^2 .
 \label{eq:maskedmomentum}
\end{equation}
The forward map stores singleton beam vectors with their positive physical
energies, while non-singleton masks are reconstructed through equivalent
relations.  This detail can change a vector's sign but not the invariant
relations used by the map.  In the inverse map, \code{fill\_momentum\_array}
implements Eq.~\eqref{eq:maskedmomentum} explicitly.

Momentum conservation implies $M_A^2=M_{\bar A}^2$.  The initialization stores
the mass of each singleton and its complement; the generator fills other masks
as they become known.  This representation makes a $t$-channel invariant just
as cheap to access as a final-state invariant mass.

\subsection{Canonical colour order and the two sets}

Let the supplied cyclic order be rotated such that label 1 is first.  If label
2 occurs at position $r$, define
\begin{equation}
 \mc A=(o_2,\ldots,o_{r-1}),\qquad
 \mc B=(o_{r+1},\ldots,o_n).
 \label{eq:sets}
\end{equation}
The masks of these lists are \code{sets(0,1)} and \code{sets(0,2)}; their
ordered entries are retained in \code{sets(1:,1)} and \code{sets(1:,2)}.
They are the final-state legs on the two arcs between the incoming legs in the
cyclic order.  The generator first connects the two arcs centrally and then
peels legs from each arc in its stored order.

This is why changing the colour order changes the sampling density even when
the final-state particle list is unchanged.  It is also why the inverse map
can evaluate a partner multichannel density only after permuting momenta into
the partner's phase-space order.

\section{Global flow of one phase-space point}

For reference, the forward map follows this sequence.

\begin{enumerate}
\item Initialise singleton masses, masks, cut-derived bounds, the canonical
  colour order, and the two sets $\mc A$, $\mc B$.
\item If PDFs are enabled, sample $\shat$ and the longitudinal centre-of-mass
  rapidity $y_{\rm cm}$.  Otherwise set $\shat=s_{\rm had}$ supplied to
  \code{init}.
\item Put the two incoming momenta back-to-back in the partonic centre-of-mass
  frame.
\item Make a central connection between $\mc A$ and $\mc B$ using either a
  two-body $t$ block or a special double-$t$ block.
\item Expand each multi-particle set.  Ordinary two-body blocks initialise a
  chain; when \code{t\_channel=.false.}, the subsequent internal steps use the
  invariant $2\to3$ kernel.
\item Divide by the global $(2\pi)^{3N-4}$ and $2\shat$ factors, then boost
  the final and incoming vectors to the lab if PDFs were used.
\end{enumerate}

The source marks an unusable point by making \code{ps\%jac} negative.  Such a
point has zero integrand and must not be interpreted as a negative physical
Jacobian.

\section{Maps from unit random numbers to invariants}

\subsection{Power maps}

For a positive variable $v\in[v_{\min},v_{\max}]$, the routine
\code{random\_to\_var} maps $x\in[0,1]$ to a density proportional to
$v^\alpha$.  For $\alpha\ne-1$ it uses
\begin{align}
 v(x)&=\left[v_{\min}^{\alpha+1}(1-x)+v_{\max}^{\alpha+1}x\right]^{1/(\alpha+1)},
 \label{eq:powermap}\\
 \left|\frac{\dd v}{\dd x}\right|&=
 \frac{v_{\max}^{\alpha+1}-v_{\min}^{\alpha+1}}{\alpha+1}\,v^{-\alpha}.
 \label{eq:powerjac}
\end{align}
For $\alpha=-1$ it instead uses the logarithmic map
\begin{equation}
 v(x)=v_{\min}^{1-x}v_{\max}^{x},\qquad
 \left|\frac{\dd v}{\dd x}\right|=v\log\frac{v_{\max}}{v_{\min}}.
 \label{eq:logmap}
\end{equation}
The derivative is multiplied into \code{ps\%jac}.  If the entire interval is
negative, the source applies Eqs.~\eqref{eq:powermap}--\eqref{eq:logmap} to
$-v$ and restores the sign afterwards.  A non-flat map across zero is rejected
and replaced by a flat map.

The default exponents are
\begin{center}
\begin{tabular}{lll}
\toprule
source variable & exponent & reason for the choice\\
\midrule
ordinary $t$ and local $s$ & \code{ip=-1} & logarithmic enhancement near a pole\\
partonic $\shat$ & \code{ip\_shat=-1.2} & stronger low-$\shat$ emphasis\\
double-$t$ invariants & \code{ip\_dt=-1} & logarithmic $t$ sampling\\
composite masses & \code{ip\_mass=-0.5} & threshold-oriented sampling\\
angles & $0$ & flat sampling\\
\bottomrule
\end{tabular}
\end{center}
Passing \code{flat=.true.} sets all four non-angular exponents to zero.  The
unused special exponent \code{101} implements
$v=v_{\min}+(v_{\max}-v_{\min})(1-\cos\pi x)/2$.

\subsection{PDF variables and the lab boost}

Let $S=s_{\rm had}=\code{this\%sqrts**2}$ and
\begin{equation}
 \tau=\frac{\shat}{S},\qquad
 x_a=\sqrt\tau\,e^{y_{\rm cm}},\qquad
 x_b=\sqrt\tau\,e^{-y_{\rm cm}} .
 \label{eq:tauy}
\end{equation}
The Jacobian is particularly simple:
\begin{equation}
 \left|\frac{\partial(x_a,x_b)}{\partial(\tau,y_{\rm cm})}\right|=1,
 \qquad \dd x_a\dd x_b=\frac{\dd\shat}{S}\dd y_{\rm cm} .
 \label{eq:pdfjac}
\end{equation}
For fixed $\tau$, the constraints $0<x_a,x_b\le1$ give
\begin{equation}
 \frac12\log\tau\le y_{\rm cm}\le-\frac12\log\tau .
 \label{eq:ylimits}
\end{equation}
The source samples $\shat$ between its cut-derived lower limit and $S$,
multiplies its Jacobian by $1/S$, then samples Eq.~\eqref{eq:ylimits} flatly.
The partonic CM vectors are transformed to the lab by a $z$ boost of rapidity
$-y_{\rm cm}$ in the source's sign convention.

\section{Composite-mass bounds and ordinary two-body blocks}

\subsection{Mass recursion}

Suppose a parent system of invariant mass $M_P$ is split into groups $I$ and
$R$.  Before cut-specific refinements, the allowed mass of $I$ is
\begin{equation}
 \left(\sum_{r\in I}m_r\right)^2
 \le M_I^2\le
 \left(M_P-\sqrt{\max(M_R^2,M_{R,\min}^2)}\right)^2 .
 \label{eq:massbounds}
\end{equation}
This is \code{shatminmax}.  \code{generate\_mass} then intersects the result
with \code{invm\_min(I)} and \code{invm\_max(I)} and samples the surviving
interval using $\alpha=\code{ip\_mass}$.  When both daughter groups are
composite, the code samples one mass first and makes the other conditional on
it.  This avoids treating two correlated masses as independent coordinates.

\subsection{$s$-channel two-body block}

For a parent $P$ with $P^2=S_P$ and daughters of masses $m_I,m_R$, work in the
rest frame of $P$.  The standard two-body kinematics are
\begin{align}
 E_I&=\frac{S_P+m_I^2-m_R^2}{2\sqrt{S_P}}, &
 E_R&=\frac{S_P-m_I^2+m_R^2}{2\sqrt{S_P}},\\
 |\boldsymbol p_I|&=|\boldsymbol p_R|=
 \frac{\sqrt{\lam(S_P,m_I^2,m_R^2)}}{2\sqrt{S_P}} .
 \label{eq:twobodykinematics}
\end{align}
Sampling $\cos\theta\in[-1,1]$ and $\phi\in[0,2\pi)$ gives the reduced
Jacobian
\begin{equation}
 J_s=\frac{\sqrt{\lam(S_P,m_I^2,m_R^2)}}{8S_P}.
 \label{eq:sblockjac}
\end{equation}
Indeed, $J_s/(2\pi)^2$ reproduces
$\dd\Phi_2=\sqrt{\lam}/(32\pi^2S_P)\,\dd\cos\theta\dd\phi$.
The source routine \code{gens\_one\_step} implements
Eqs.~\eqref{eq:twobodykinematics}--\eqref{eq:sblockjac} through
\code{mom2cx} followed by \code{boostm}.

\subsection{$t$-channel two-body block}

The physical range of a momentum transfer follows from an ordinary two-body
scattering problem.  Let $X$ be the squared total energy, let incoming squared
masses be $W,Z$, let outgoing squared masses be $U,V$, and define the transfer
to the outgoing object of mass $U$ by $T=(p_{\rm in}-p_U)^2$.  At fixed $X$,
the endpoints are
\begin{equation}
 T_{1,2}=U+W-
 \frac{(X+U-V)(X+W-Z)\mp
 \sqrt{\lam(X,U,V)\lam(X,W,Z)}}{2X}.
 \label{eq:tlimitsgeneric}
\end{equation}
The order of the two signs is immaterial: the code explicitly takes
\code{min} and \code{max}.  This is the formula in \code{tminmax}, where
\begin{equation}
 (X,Z,U,V,W)=
 (\hat s_i,t_i,\hat s_{i-1},m_i^2,m_a^2).
 \label{eq:tminmaxarguments}
\end{equation}

After clipping Eq.~\eqref{eq:tlimitsgeneric} by cuts, \code{gent\_one\_step}
samples the relevant transfer and an azimuth.  In a suitable parent rest frame,
the transfer fixes the longitudinal component of a daughter:
\begin{equation}
 p_{I,z}=-\frac{m_I^2+m_a^2-T-2E_I E_a}{2|\boldsymbol p_a|},
 \qquad p_{I,\perp}^2=|\boldsymbol p_I|^2-p_{I,z}^2 .
 \label{eq:treconstruct}
\end{equation}
The source rotates the vector back to the actual incoming direction and boosts
it to the event frame.  Its reduced Jacobian is accumulated exactly as
\begin{equation}
 J_t^{\rm code}=
 \frac{1}{4\sqrt{\lam\bigl(M_{I+R}^2,0,M_{I+R+b}^2\bigr)}}.
 \label{eq:tblockjac}
\end{equation}
The apparently unusual third argument in Eq.~\eqref{eq:tblockjac} is a
consequence of the crossed mask convention.  It should not be replaced by a
textbook $2\to2$ factor without following the recursive invariant definitions.

When the parent is the full final-state system, the code additionally uses
\begin{equation}
 E_I=\frac{\sqrt{\shat}+(M_I^2-M_R^2)/\sqrt{\shat}}{2},\qquad
 E_R=\frac{\sqrt{\shat}+(M_R^2-M_I^2)/\sqrt{\shat}}{2}
 \label{eq:twobodyenergies}
\end{equation}
to restrict sampled masses by the precomputed $\ET$ minima.  These are useful
necessary conditions, not a replacement for the final cut test.

\section{The special double-$t$ block}

\code{double\_t(i,ir,ia,ib)} is used when one colour-order set consists of a
single leg $i$ and the other is a multi-particle remainder $r$.  It generates
three variables $t_a,t_b,\phi$ while deriving the remainder mass.  In the
partonic CM frame let
\begin{equation}
 t_a=(p_i-p_a)^2,\qquad t_b=(p_i-p_b)^2,
 \qquad p_a^2=p_b^2=0 .
 \label{eq:doubletdefs}
\end{equation}
Solving these equations gives precisely the source formulas
\begin{align}
 E_i&=\frac{-t_a-t_b+2m_i^2}{2\sqrt{\shat}}, &
 p_{i,z}&=\frac{t_a-t_b}{2\sqrt{\shat}},\\
 p_{i,\perp}^2&=
 \frac{t_at_b+m_i^4-m_i^2(t_a+t_b)}{\shat}-m_i^2,
 & M_r^2&=\shat+t_a+t_b-m_i^2 .
 \label{eq:doubletsolution}
\end{align}
The azimuth determines the transverse components
$p_{i,x}=p_{i,\perp}\cos\phi$ and
$p_{i,y}=p_{i,\perp}\sin\phi$; the remainder is $p_r=p_a+p_b-p_i$.

If the remainder must have $M_r^2\geq M_{r,\min}^2$, the initial range is
\begin{equation}
 t_a\in\left[
 \frac{-\shat+m_i^2+M_{r,\min}^2-
 \sqrt{\lam(\shat,m_i^2,M_{r,\min}^2)}}{2},
 \frac{-\shat+m_i^2+M_{r,\min}^2+
 \sqrt{\lam(\shat,m_i^2,M_{r,\min}^2)}}{2}\right].
 \label{eq:doubletfirstbound}
\end{equation}
After $t_a$ is chosen, $M_r^2\geq M_{r,\min}^2$ and
$p_{i,\perp}^2\geq0$ imply, respectively,
\begin{align}
 t_b&\geq-\shat-t_a+m_i^2+M_{r,\min}^2,\\
 t_b&\leq \frac{m_i^2(m_i^2-\shat-t_a)}{m_i^2-t_a} .
 \label{eq:doubletsecondbound}
\end{align}
The routine further intersects both intervals with mask-level cut bounds and
with its $\ET$ restrictions.  Its reduced local factor is
\begin{equation}
 J_{tt}=\frac{1}{4\sqrt{\lam(M_{i+r}^2,0,0)}}=
 \frac{1}{4\shat}
 \quad\hbox{for a massless incoming pair.}
 \label{eq:doubletjac}
\end{equation}

\section{The invariant $2\to3$ Byckling--Kajantie block}

\subsection{Local invariant dictionary}

The routine \code{gen23\_one\_step(i,ir,ib,im1)} is easiest to read by
renaming its masks locally.  Let $q=\code{i}$ be the leg being resolved,
$r=\code{ir}$ the remaining group, $h=\code{im1}$ the neighbouring already
known leg, and $b=\code{ib}$ the beam-side mask.  Define
\begin{align}
 H_0&=(p_q+p_r)^2,& H_-&=p_r^2,& H_+&=(p_h+p_q+p_r)^2,\\
 T_0&=(p_q+p_r+p_b)^2,& T_-&=(p_r+p_b)^2,&
 T_+&=(p_h+p_q+p_r+p_b)^2,\\
 s&=(p_q+p_h)^2,& m_q^2&=p_q^2,&m_h^2&=p_h^2 .
 \label{eq:bkdictionary}
\end{align}
For example, the code call
\begin{center}
\code{sminmax(invm(ir+i), invm(ir), invm(ir+i+im1),\ldots)}
\end{center}
passes $(H_0,H_-,H_+,T_0,T_-,T_+,m_q^2,m_h^2)$ in that order.  This local
notation avoids an otherwise easy off-by-one ambiguity in the historical
$\hat s_i,t_i,s_i$ labels.

If $q$ or $r$ is composite, its invariant mass is generated first using
Eq.~\eqref{eq:massbounds}.  The continuous invariants particular to this
local $2\to3$ relation are then $T_-$ and $s$.  The remaining azimuth is not
sampled independently: $s$ determines $\cos\phi$, leaving two possible
orientations.

\subsection{Gram determinant and local measure}

The source constructs the symmetric matrix
\begin{equation}
 A_4=\begin{pmatrix}
 0&T_--H_-&T_0-H_0&T_+-H_+\\
 T_--H_-&2T_-&T_0+T_--m_q^2&T_-+T_+-s\\
 T_0-H_0&T_0+T_--m_q^2&2T_0&T_0+T_+-m_h^2\\
 T_+-H_+&T_-+T_+-s&T_0+T_+-m_h^2&2T_+
 \end{pmatrix},\qquad
 \Delta_4=\frac{\det A_4}{16}.
 \label{eq:gram4}
\end{equation}
It is the quantity returned by \code{gram\_determinant4}.  In the physical
interior relevant here, $\Delta_4<0$.  After the masses are fixed, the
invariant reduced measure used by the code is
\begin{equation}
 \dd\overline\Phi_{23}=\frac{\dd T_-\,\dd s}{8\sqrt{-\Delta_4}}.
 \label{eq:bkmeasure}
\end{equation}
The factor in Eq.~\eqref{eq:bkmeasure} is accumulated at the end of
\code{gen23\_one\_step}; the power-map derivatives for $T_-$ and $s$ have
already been accumulated separately.

\subsection{First boundary: the transfer $T_-$}

Applying Eq.~\eqref{eq:tlimitsgeneric} to the local dictionary gives
\begin{align}
 T_-^{(1,2)}={}&H_- -
 \frac{(H_0+H_--m_q^2)(H_0-T_0)
 \mp\sqrt{\lam(H_0,H_-,m_q^2)\lam(H_0,0,T_0)}}{2H_0},\\
 T_-&\in[\min(T_-^{(1)},T_-^{(2)}),\max(T_-^{(1)},T_-^{(2)})].
 \label{eq:bktbound}
\end{align}
Here the incoming mass argument is zero, as required by the implementation.
The source clips this interval with \code{invm\_min(ir+ib)} and
\code{invm\_max(ir+ib)}, then with bounds derived from transverse-energy
requirements.  Empty intervals are rejected before a random number is mapped.

For the latter refinement, boost to a frame in which $p_q+p_r$ has zero
longitudinal momentum.  If $E_{qr}$ and $E_b$ are the corresponding energies
and $\mu_q,\mu_r$ are the effective transverse-energy minima, define
\begin{equation}
 B=E_{qr}^2-\mu_q^2+\mu_r^2,
 \qquad R=\lam(E_{qr}^2,\mu_q^2,\mu_r^2).
 \label{eq:etbounddefinitions}
\end{equation}
The code intersects the transfer interval with
\begin{equation}
 T_-\geq m_r^2-\frac{E_b}{E_{qr}}(B+\sqrt R),\qquad
 T_-\leq m_r^2-\frac{E_b}{E_{qr}}(B-\sqrt R).
 \label{eq:ettransferbound}
\end{equation}
The effective $\mu$ values include a correction when a singleton daughter
already carries transverse momentum.  This makes Eq.~\eqref{eq:ettransferbound}
stronger than merely applying $\ET\geq\ET^{\min}$ after reconstruction.

\subsection{Second boundary: the adjacent invariant $s$}

The $s$ boundary is obtained by demanding $|\cos\phi|\leq1$.  The source
first defines
\begin{equation}
 V=-\frac{1}{8}\det\begin{pmatrix}
 2H_0&H_0-T_0&H_0+H_+-m_h^2\\
 H_0-T_0&0&H_+-T_+\\
 H_0+H_--m_q^2&H_--T_-&0
 \end{pmatrix}.
 \label{eq:vdet}
\end{equation}
The matrix in Eq.~\eqref{eq:vdet} is the code's $V$ matrix; it is not itself a
symmetric Gram matrix.  Define the polynomial
\begin{align}
 G(x,y,z,u,v,w)={}&x^2y+xy^2+z^2u+zu^2+v^2w+vw^2
 +xzw+xuv+yzv+yuw\notag\\
 &{}-xy(z+u+v+w)-zu(x+y+v+w)-vw(x+y+z+u),
 \label{eq:gpoly}
\end{align}
and
\begin{equation}
 \Gamma=G(T_0,H_+,H_0,T_+,m_h^2,0)\,
          G(T_-,H_0,H_-,T_0,m_q^2,0).
 \label{eq:gamma}
\end{equation}
The exact kinematic endpoints are
\begin{equation}
 s_\pm=H_-+H_++
 \frac{2\left(4V\pm\sqrt\Gamma\right)}{\lam(H_0,T_0,0)},
 \qquad s\in[\min(s_+,s_-),\max(s_+,s_-)].
 \label{eq:sbounds}
\end{equation}
These equations are implemented by \code{computeV}, \code{G}, and
\code{sminmax}.  The source subsequently intersects Eq.~\eqref{eq:sbounds}
with direct mask bounds for $q+h$ and, when $h$ is a final-state leg, with
additional $\ET$ and neighbouring-$\Delta R$ pre-bounds.

For completeness, the relation inverted by \code{getphifroms} is
\begin{equation}
 \cos\phi=
 \frac{\tfrac12(s-H_--H_+)\lam(H_0,T_0,0)-4V}{\sqrt\Gamma}.
 \label{eq:cosphi}
\end{equation}
Equation~\eqref{eq:sbounds} is simply Eq.~\eqref{eq:cosphi} evaluated at
$\cos\phi=\pm1$.  This is the analytic origin of both the $s$ limits and the
two-fold ambiguity.

\subsection{Momentum reconstruction and the two branches}

Given $T_-$ and a solution for $\phi$, \code{gentcms2} reconstructs the two
outgoing objects in the rest frame of a suitable parent.  It aligns the
massless reference momentum with the $z$ axis, uses the transfer to fix the
longitudinal component, uses the positive square root for the transverse
magnitude, and fixes the azimuth relative to the known reference vector $h$.
It then rotates and boosts back.  This construction enforces all invariants in
Eq.~\eqref{eq:bkdictionary}.

For an interior value of $s$, Eq.~\eqref{eq:cosphi} has
\begin{equation}
 \phi_1=\arccos(\cos\phi),\qquad
 \phi_2=2\pi-\arccos(\cos\phi).
 \label{eq:twobranches}
\end{equation}
The base factor in Eq.~\eqref{eq:bkmeasure} accounts for both branches.  If
both reconstructed points satisfy the local $\ET$ requirements, an additional
flat random number selects one with probability $1/2$.  If exactly one branch
survives, the source divides the Jacobian by two.  If neither survives, the
forward map rejects the point.  The extra random number does not describe a
new physical dimension; it only resolves Eq.~\eqref{eq:twobranches}.

\section{Recursive topology and number of extra random numbers}

Let $n_A=|\mc A|$ and $n_B=|\mc B|$.  The central operation is chosen as
follows.

\begin{center}
\begin{tabular}{lll}
\toprule
condition & central routine & interpretation\\
\midrule
$n_A>1$, $n_B>1$ & \code{gent\_one\_step} & central $t$-channel block\\
$n_A=1$, $n_B>1$ & \code{double\_t} & singleton against a composite arc\\
$n_A>1$, $n_B=1$ & \code{double\_t} & reflected singleton case\\
$n_A=n_B=1$ & \code{gent\_one\_step} & ordinary $2\to2$ scattering\\
\bottomrule
\end{tabular}
\end{center}

Each set with three or more legs is then opened with an ordinary $t$ block and
continued one leg at a time.  For later steps the choice is
\begin{equation}
 \code{t\_channel=.true.}:\quad\code{gent\_one\_step},
 \qquad
 \code{t\_channel=.false.}:\quad\code{gen23\_one\_step}.
 \label{eq:channelchoice}
\end{equation}
For a two-leg set opposite a non-empty set, the same choice applies once.  A
two-leg set with an empty opposite side uses the ordinary two-body block.

Each invariant $2\to3$ step needs one branch selector.  The source therefore
allocates
\begin{align}
 n_{\rm extra}={}&\boldsymbol1_{n_A=2,\,n_B>0}
 +\boldsymbol1_{n_B=2,\,n_A>0}\notag\\
 &{}+\max(n_A-2,0)+\max(n_B-2,0).
 \label{eq:extra}
\end{align}
These coordinates appear after the adaptive coordinates in \code{ps\%x}.
They are intentionally sampled flatly by the integrator.  As discussed in the
audit, the formula is not conditional on \code{t\_channel}; it can therefore
allocate unused selectors in pure-$t$ mode.

\section{Cut-derived bounds}

\subsection{Cuts ultimately tested}

The generator is only a proposal density.  The final acceptance function in
\code{cuts.f03} checks, for each physical external momentum,
\begin{equation}
 p_{\mathrm{T},i}\geq p_{\mathrm{T},i}^{\rm cut},\qquad |\eta_i|\leq\eta_i^{\rm max},
 \qquad |(p_i+p_j)^2|\geq Q_{ij}^2,
 \qquad \Delta R_{ij}\geq R_{ij}^{\rm cut}.
 \label{eq:finalcuts}
\end{equation}
The phase-space initialization tries to turn a subset of these into integration
limits.  This can improve efficiency dramatically, but a pre-bound must be a
\emph{necessary} condition for the final cuts.  A bound that is stronger than
necessary removes valid phase space and biases an integral.

For a final-only mask $A$, the code constructs
\begin{equation}
 c_{ij}^{\rm code}=\max\left[Q_{ij}^2,
 2p_{\mathrm{T},i}^{\rm cut}p_{\mathrm{T},j}^{\rm cut}
 \left(1-\cos R_{ij}^{\rm cut}\right)\right],
 \qquad C_A=\sum_{i<j\in A}c_{ij}^{\rm code}.
 \label{eq:codepairbound}
\end{equation}
It then sets
\begin{equation}
 M_{A,\min}^{2\,\rm code}=
 \begin{cases}
 \max\left[(\sum_{i\in A}m_i)^2,\ \dfrac{|A|}{|A|-1}C_A\right],
      &A=\mc F\quad\hbox{(all final legs)},\\[1.1ex]
 \max\left[(\sum_{i\in A}m_i)^2,\ \dfrac12 C_A\right],
      &A\subsetneq\mc F.
 \end{cases}
 \label{eq:codeaggregatebound}
\end{equation}
This is \code{setup\_PS\_cuts}, lines 248--262.  The corresponding transverse
energy proposal threshold is
\begin{equation}
 E_{\mathrm{T},A}^{\min,\rm code}=\max\left[\sum_{i\in A}\sqrt{m_i^2+(p_{\mathrm{T},i}^{\rm cut})^2},\
 \sqrt{M_{A,\min}^{2\,\rm code}}\right].
 \label{eq:etmin}
\end{equation}
For the full final state, $\shat$ is sampled above
\begin{equation}
 \hat{s}_{\min}^{\rm code}=\max\left[M_{\mc F,\min}^{2\,\rm code},
 (E_{\mathrm{T},\mc F}^{\min,\rm code})^2\right].
 \label{eq:shatmincode}
\end{equation}

For a two-particle mask containing one incoming and one final leg, the source
uses a negative $t$ upper bound,
\begin{equation}
 t_{if}^{\max,\rm code}=-\max\left(Q_{if}^2,
 (p_{\mathrm{T},i}^{\rm cut})^2,(p_{\mathrm{T},f}^{\rm cut})^2\right),
 \label{eq:tcutcode}
\end{equation}
and assigns the same bound to the complementary mask.  This is appropriate as
a massless crossed-pair relation, subject to the caveats in
Section~\ref{sec:audit}.

\subsection{Exact identities useful for checking bounds}

For a final-only set $A$ of $r$ particles,
\begin{equation}
 M_A^2=\sum_{i<j\in A}s_{ij}-(r-2)\sum_{i\in A}m_i^2,
 \qquad s_{ij}=(p_i+p_j)^2.
 \label{eq:subsetidentity}
\end{equation}
Consequently, if $s_{ij}\geq L_{ij}$ are genuine pairwise lower bounds, a
safe aggregate lower bound is
\begin{equation}
 M_A^2\geq
 \max\left\{(\sum_{i\in A}m_i)^2,
 \sum_{i<j\in A}L_{ij}-(r-2)\sum_{i\in A}m_i^2\right\}.
 \label{eq:exactaggregatebound}
\end{equation}
For massless legs this reduces to $M_A^2\geq\sum_{i<j}L_{ij}$.

For massless final legs with $p_{\mathrm{T},i}\geq p_i$ and a usual
$0\leq R_{ij}\leq\pi$ cut,
\begin{equation}
 s_{ij}=2p_{\mathrm{T},i}p_{\mathrm{T},j}\left(\cosh\Delta\eta_{ij}-\cos\Delta\phi_{ij}\right)
 \geq2p_ip_j(1-\cos R_{ij}).
 \label{eq:drbound}
\end{equation}
This establishes the massless part of Eq.~\eqref{eq:codepairbound}.  It is a
lower bound, not generally the exact minimum after all other cuts and momentum
conservation are imposed.

The remaining implementation refinements use these ideas locally: parent
masses are restricted by Eq.~\eqref{eq:massbounds}, $t$ intervals by
Eqs.~\eqref{eq:bktbound} and \eqref{eq:ettransferbound}, and $s$ intervals by
Eq.~\eqref{eq:sbounds} plus the mask bounds.  Any cuts not encoded in these
limits are left to Eq.~\eqref{eq:finalcuts}.

\section{The inverse map}

\code{gen23\_compute\_x\_from\_momenta} is used by the colour-singlet
multichannel weight computation.  Given a physical event, it:

\begin{enumerate}
\item reconstructs every mask vector with Eq.~\eqref{eq:maskedmomentum} and
  boosts them to the partonic CM frame;
\item obtains $\shat$, $\tau$, and $y_{\rm cm}$ if PDFs are enabled;
\item follows the same central and recursive topology as the forward map;
\item evaluates each generated invariant from the supplied momenta and applies
  \code{var\_to\_random}, the analytic inverse of
  \code{random\_to\_var};
\item recomputes the same forward Jacobian.  This is the quantity needed to
  compare proposal densities across channels.
\end{enumerate}

The inverse uses the same $t$, $s$, mass, and $\ET$ bounds, which is essential:
a momentum outside a partner channel's support should return a negative
Jacobian and be omitted from that partner's multichannel density.  The branch
coordinate itself is deliberately not reconstructed, since it is an auxiliary
flat selector rather than an adaptive physical coordinate.

\section{Implementation audit and improvements}
\label{sec:audit}

This section distinguishes exact issues from deliberate but loose bounds.  The
findings below are based on the present source and on the identities derived
above; no code change is made by this document.

\subsection{Confirmed support or semantic discrepancies}

\paragraph{All-final aggregate bound can exclude accepted phase space
(high priority).}
For massless final legs, Eq.~\eqref{eq:subsetidentity} says that a full
final-state system subject to pair lower bounds has the necessary condition
$\shat\geq C_{\mc F}$, not
$\shat\geq N C_{\mc F}/(N-1)$.  The latter is the first branch of
Eq.~\eqref{eq:codeaggregatebound}.  A particularly transparent counterexample
is $N=2$, zero masses and zero $\pt$ cuts, with only
$s_{34}\geq Q^2$.  The final cut accepts every $\shat=s_{34}\geq Q^2$.
The current setup instead sets
\begin{equation}
 \code{invm\_min(all finals)}=2Q^2,
\end{equation}
so PDF sampling starts at $2Q^2$ and the valid interval
$Q^2\leq\shat<2Q^2$ is never proposed.  The same issue exists for more legs
unless another valid bound happens to dominate.  This is a coverage issue, not
just an efficiency issue.

The direct replacement should be a demonstrably necessary bound such as
Eq.~\eqref{eq:exactaggregatebound}; for massless legs, use $C_A$ rather than
$|A|C_A/(|A|-1)$.  A regression test should integrate a constant matrix element
with a non-zero pair-mass cut and compare against an independent phase-space
volume or a deliberately loose generator.

\paragraph{Initial--final pair bounds overwrite the requested pair cut
(efficiency and regulator risk).}
For a two-leg mask with one beam and one final leg,
\code{setup\_PS\_cuts} loops over all four ordered pairs including both
diagonal pairs (source lines 232--239).  It assigns instead of accumulates
\code{invm\_max(k)}.  The final diagonal iteration normally has
\code{sqrt\_s\_min(i,i)=0}, so an off-diagonal $Q_{if}$ can be overwritten by
the $\pt$ term.  The later final cut still rejects forbidden points, so this
does not by itself lose physical support; it can however remove the intended
$t$-channel pre-regulation and severely reduce efficiency.

The fix is to loop only over the two distinct legs and retain the strongest
negative upper bound, e.g. initialise to zero and apply
\code{invm\_max(k)=min(invm\_max(k),-candidate)}.  Symmetrise or validate the
input cut matrix at initialization.

\paragraph{Massive initial--final pair semantics differ from the final cut
(conditional coverage risk).}
For a massless beam leg $a$ and a massive final leg $f$,
\begin{equation}
 s_{af}^{(+)}=(p_a+p_f)^2=m_f^2+2p_a\!\cdot p_f,
 \qquad t_{af}=(p_a-p_f)^2=2m_f^2-s_{af}^{(+)} .
 \label{eq:massivetrelation}
\end{equation}
The final cut in Eq.~\eqref{eq:finalcuts} tests $s_{af}^{(+)}$, whereas
Eq.~\eqref{eq:tcutcode} implements $t_{af}\leq-Q_{af}^2$.  If the API intends
these to represent the same $Q_{af}$ cut, the latter requires
$s_{af}^{(+)}\geq Q_{af}^2+2m_f^2$, which is stricter than the final test.
For massless legs the two expressions coincide.  The implementation should
either define incoming--final cuts explicitly as crossed $-t$ cuts or derive
the mask bound from Eq.~\eqref{eq:massivetrelation} so that it matches the
final acceptance definition.

\paragraph{The inverse $2\to3$ branch test is not symmetric with the forward
map (multichannel-density risk).}
The forward routine rejects when neither $\phi_1$ nor $\phi_2$ satisfies the
local $\ET$ checks (lines 795--831).  The inverse routine tests ``both'',
``only first'', and ``only second'' (lines 1677--1683), but has no final
\code{else} that sets \code{ps\%jac} negative.  Thus a point for which neither
candidate satisfies those local conditions can continue with a non-zero
partner density.  This matters precisely in the multichannel use case, where a
point generated in one channel is tested against another channel's support.

The inverse should mirror the forward rejection branch.  Add an explicit
``neither'' case that sets a negative Jacobian and returns, then test
forward--inverse agreement for randomly chosen partner channel permutations.

\paragraph{The interface calls an input a rapidity cut but \code{gen23} does
not use it, and the later cut is pseudorapidity.}
\code{gen23\_init} explicitly labels \code{rap\_cut} ``not used'' and does
not store it.  The later acceptance routine applies $|\eta|<\eta_{\max}$,
not true rapidity $|y|<y_{\max}$.  For massless particles $y=\eta$, but for
massive legs they differ.  Therefore neither the name nor the phase-space
pre-bound semantics are currently those of a massive-particle rapidity cut.
Rename the input to an eta cut, or implement the desired rapidity definition
consistently and, where possible, include it in the proposal bounds.

\subsection{Bounds that are safe but unnecessarily loose}

For a proper final-state subset, the code uses $C_A/2$ in
Eq.~\eqref{eq:codeaggregatebound}.  Since Eq.~\eqref{eq:exactaggregatebound}
is available, this can be much weaker than necessary.  It does not bias the
result, but it sends avoidable points into later rejection tests.  The same is
true of the massless $\Delta R$ lower estimate when legs are massive or when
other kinematic information is known.

A stronger safe massless pair estimate for $0\le R\le\pi$ is already
Eq.~\eqref{eq:drbound}.  For massive legs, the exact relation depends on the
rapidity rather than the pseudorapidity used by \code{cuts.f03}; it should not
be upgraded casually.  A robust implementation can first use the exact
mass-threshold and pair-mass identity in Eq.~\eqref{eq:exactaggregatebound},
then add only separately proved $\pt$/$\Delta R$ inequalities.

\subsection{Numerical robustness issues}

\begin{itemize}
\item The logarithmic map needs strictly positive endpoints.  With fully
  inclusive massless $t$ ranges, an endpoint can be zero, but
  \code{ip=-1} evaluates logarithms or powers of zero.  Likewise
  \code{ip\_shat=-1.2} is not defined at a zero lower $\shat$ endpoint.
  Require a positive physical regulator, switch to a finite endpoint-aware
  map, or return a clean rejected point before evaluating the map.
\item \code{getphifroms} returns $\phi=0$ when the computed cosine lies
  outside $[-1,1]$, even if the excess is not roundoff.  That does not preserve
  the requested invariant $s$.  Use a scale-aware tolerance: clamp only a
  tiny overshoot and reject a material violation.
\item Several expressions take square roots of quantities that can be a few
  ulps negative near an endpoint, and several azimuths use
  \code{atan(y/x)} rather than \code{atan2(y,x)}.  Use an invariant-scale
  tolerance, clamp only tiny negative radicands, and use \code{atan2} to
  avoid divisions by zero and quadrant errors.
\item The Gram and $V$ determinants use an absolute LU pivot tolerance
  \code{1d-8}.  Their entries carry powers of energy, so this criterion is
  not scale invariant.  Rescaling the matrix or using a relative pivot test
  would make high- and low-energy behaviour more uniform.
\item The global module variables \code{includePDF}, \code{ip},
  \code{ip\_shat}, \code{ip\_dt}, and \code{ip\_mass} are shared by all
  \code{phase\_space\_gen23} objects.  This is safe only if all active maps
  use identical settings and calls are serialized.  Put them in \code{this}
  for re-entrant or threaded use.
\end{itemize}

\subsection{Topology-selection and bookkeeping caveats}

The public comment says \code{t\_chan=.false.} selects the Byckling--Kajantie
construction and \code{.true.} selects the ordinary $t$ construction.  In
practice the central split between two non-trivial sets is controlled by the
private compile-time parameter \code{use\_t\_channel\_at\_start=.true.}.
Thus even \code{t\_chan=.false.} begins with a $t$ block in that topology.
This is a valid sampling choice, but the public wording should describe the
mixed construction accurately.

Also, Eq.~\eqref{eq:extra} allocates branch-selector dimensions independently
of \code{t\_channel}.  In pure-$t$ mode those coordinates are unused.  This
does not change a properly normalised flat integral, but it adds dimensions to
the integrator and can slow adaptation.

Finally, the nested Breit--Wigner routines are present but their calls are
commented out in the forward map.  The current \code{gen23} implementation
does \emph{not} perform resonance-mass importance sampling through those
routines.

\section{Recommended validation programme}

The following tests exercise the derivation and catch the issues above.

\begin{enumerate}
\item \textbf{Constant-integrand volumes.}  Compare $2\to2$ and $2\to3$
  massless volumes against analytic results, both with flat maps and with the
  default importance maps.
\item \textbf{Threshold scan.}  With PDFs enabled, set a single final--final
  $Q_{ij}$ cut and verify non-zero support immediately above $Q_{ij}^2$.
  This directly tests the all-final bound in Eq.~\eqref{eq:codeaggregatebound}.
\item \textbf{Forward/inverse closure.}  Generate points in every relevant
  topology, invert them in the same channel, regenerate, and compare all
  momenta, invariants, $x$ coordinates, and Jacobians within a scale-aware
  tolerance.
\item \textbf{Partner-channel support.}  Permute a generated point into each
  partner colour-order channel.  Check that an inverse map returns negative
  Jacobian exactly when neither azimuthal branch is supported.
\item \textbf{Boundary stress tests.}  Probe $\pt\to0$, $t\to0$, Gram
  determinant endpoints, massive legs, $\Delta R$ limits, and nearly
  threshold composite masses with floating-point exceptions enabled.
\end{enumerate}

\section{Summary}

\code{gen23} is a hybrid recursive phase-space generator.  Its invariant
$2\to3$ core samples a transfer and an adjacent pair mass, derives their
physical region through Gram determinants, reconstructs one of two azimuthal
branches, and uses the factor $1/(8\sqrt{-\Delta_4})$.  Ordinary two-body and
double-$t$ blocks complete the colour-order topology.  The total event weight
is the product of all coordinate derivatives and local reduced Jacobians,
divided by $(2\pi)^{3N-4}$ and $2\shat$.

The analytic formulae also make the implementation reviewable.  In particular,
the full-final-state pair-cut lower bound should be corrected before relying on
PDF integrations with non-zero pair-mass cuts; the inverse no-valid-branch
case should mirror the forward rejection; and endpoint/numerical handling
should be hardened.  The other improvements listed above mostly preserve
coverage while making the proposal density more efficient and robust.

\begin{thebibliography}{9}
\bibitem{BK}
E.~Byckling and K.~Kajantie,
``Reductions of the Phase-Space Integral in Terms of Simpler Processes,''
\emph{Physical Review} \textbf{187} (1969) 2008--2016,
\href{https://doi.org/10.1103/PhysRev.187.2008}{doi:10.1103/PhysRev.187.2008}.
\end{thebibliography}

\end{document}
